% ===== tr-08 · b15 · tr08d：第 5 章 p34--37（原文页 34--37）=====
% [原文 p.34]
关于上述问题，如果人们在有限个候选收缩映射之间难以取舍，可以过渡到一个公共的 snc 消解（或许还要过渡到有限基变换，以得到半稳定的 snc 消解）。当然，我们对模型 $\mathcal X$ 爆破得越多，比较性质就越弱。

\begin{remark}{5.13}
Pille-Schneider 与 Mazzon 近来的工作 [59] 提出把联系于若干因子化对数终端模型的收缩映射粘合起来，以得到一个映射 $X_K^{an}\to Sk(X)$。他们的映射仍然经由某个更大 snc 模型的对偶相交复形分解，因此与比较性质相容。
\end{remark}

% [Tutor Note: 注 13 的 "divides" 无关，关键是 "still factors through the dual intersection complex of some larger snc model"------即新映射仍可分解为 $X_K^{an}\to\Delta_{\mathcal X'}\to Sk(X)$，故与定义 5.5.6 的比较性质不冲突。Pille-Schneider 与 Mazzon 是两个姓氏（Pille-Schneider 为复姓），不要拆成三人。]

\begin{itemize}
\item 若没有比较性质，似乎很难对 $\|\cdot\|_{CY}$ 给出任何微分几何解释，因为 $X_K^{an}$ 含有任意大的对偶相交复形，因而高度复杂。这一假设情形的启发式直观是：Calabi--Yau 度量的势论极限将需要无穷多次爆破才能描述。这一点尚未被任何定理排除；其可信度留给读者判断。
\end{itemize}

留给代数几何学家的公开问题是

\begin{question}{5.6}
能否对足够大的一类例子（例如来自 Gross--Siebert 纲领 [29] 的那些）证明比较性质？
\end{question}

% [Tutor Note: 问题编号按原文为 Question 6，本章（第 5 章）内套用 $\{5.6\}$；与本章的式 (5.6)、注 (5.6) 不是同一计数器，互不冲突。原文 p33 的 itemize 有 4 条、p34 这条 "Without the comparison property" 是同一评论列表的第 5 条，中间被一段非列表段落（"The formulation of the comparison property..."）打断；译文按原文的分段结构在新的一段处另起一个 itemize。]

\section{证明策略一览}
我们讨论 SYZ 猜想的弱度量版本的一些近期进展。

\begin{theorem}{5.6.1}
[53] 设 $X\to S\setminus\{0\}$ 是 Calabi--Yau 流形的大复结构极限，其极化丰沛线丛为 $L\to X$。假设 $X$ 满足 NA MA 与实 MA 比较性质（参看第 5.5.6 节）。则对充分小的 $t\in S\setminus\{0\}$，在 $X_t$ 的某个开子集上存在关于 Calabi--Yau 结构 $(\omega_{CY,t},\Omega_t)$ 的特殊拉格朗日 $T^n$-纤维化，其正规化 Calabi--Yau 测度当 $t\to 0$ 时趋于 $100\%$。
\end{theorem}

% [Tutor Note: 定理编号 = 章号.原文编号：本章为第 5 章，故原文 "Theorem 6.1" 记作 $\{5.6.1\}$（不是 6.1）。同理下文 Theorem 6.2、Proposition 6.3、6.4 分别记 $\{5.6.2\}$、$\{5.6.3\}$、$\{5.6.4\}$。$S\setminus\{0\}$ 中的 $S$ 是底曲线，与后文 snc 模型记号无关。]

存在一族非常特殊的射影 Calabi--Yau 超曲面，对它可以绕过 NA MA 与实 MA 比较性质，代价是过渡到子序列：
\begin{equation}
X_t=\Big\{Z_0Z_1\dots Z_{n+1}+t\sum_{i=0}^{n+1}Z_i^{n+2}=0\Big\}\subset\C\mathrm{P}^{n+1},\qquad t\in\R,\ 0<t\ll 1. \tag{5.18}
\end{equation}
我们称之为 Fermat 族，因其与著名的 Fermat 多项式 $\sum_{i=0}^{n+1}Z_i^{n+2}$ 有关。

% [Tutor Note: (18) 中有 $n+2$ 个变量 $Z_0,\dots,Z_{n+1}$，每一项的次数都是 $n+2$，这就是射影空间取 $\PP^{n+1}$（不是 $\PP^{n+2}$）的原因；求和上下限为 $i=0$ 到 $n+1$，不要漏掉 $Z_{n+1}$。$t$ 取实数且 $0<t\ll 1$（原文 $t\in\R$，与 (15) 的复参数 $t$ 含义不同）。]

% [原文 p.35]
\begin{theorem}{5.6.2}
[52] 对 Fermat 族，考虑 $X_t$ 上属于极化类 $\frac{1}{|\log|t||}c_1(\mathcal O(1))$ 的 Calabi--Yau 度量 $\omega_{CY,t}$。则当 $t\to 0$ 时，沿 $X_t$ 的一个子序列，在 $X_t$ 的某个开子集上存在特殊拉格朗日 $T^n$-纤维化，其正规化 Calabi--Yau 测度趋于 $100\%$。
\end{theorem}

% [Tutor Note: 极化类系数原文排成竖式分数 $\frac{1}{|\log|t||}c_1(\mathcal O(1))$，分母是 $|\log|t||$（外层绝对值套住 $\log|t|$），不是 $\log|t|$ 的绝对值再取倒数以外的其他结合方式。$\mathcal O(1)$ 用花体 O（$\PP^{n+1}$ 的超平面丛）。]

我们的叙述将集中于主要思路及其诸多细微之处，而不给出证明的全部细节。

\subsection{化为势的估计}
两策略的共同部分是把特殊拉格朗日的存在性问题化为势的 $C^0$-估计。

给定固定的 snc 模型 $\mathcal X\to S$，存在对数映射 $\mathrm{Log}_{\mathcal X}:X_t\to\Delta_{\mathcal X}$，它在 $O\big(\frac{1}{|\log|t||}\big)$ 的坐标歧义下确定。给定 $Sk(X)\subset\Delta_{\mathcal X}$ 的一个 $n$ 维面 $\Delta_J$，我们考虑 $\Delta_J$ 内部略微收缩后的版本在对数映射下的原像 $U_{J,t}\subset X_t$。我们将随意地把 $U_{J,t}$ 收缩若干次，只要被删去的集合在 $t\to 0$ 极限下具有可忽略的 Calabi--Yau 测度。由于 $U_{J,t}$ 可以看作 $(\C^*)^n$ 的环面不变子集，通过过渡到带坐标 $\frac{\log z_i}{\log|t|}$ 的万有覆盖，我们可以使 $C^k_{loc}$ 范数在 $t$ 上一致地有意义。

第一个主要步骤是把 $C^0$-估计改进为 $C^\infty$-估计。

\begin{proposition}{5.6.3}[53, section 4.5]
设 $\phi_0$ 是实 MA 方程 (5.17) 在 $\Delta_J$ 内部的一个 Alexandrov 解。假设 $U_{J,t}$ 上的 Calabi--Yau 度量容许局部势函数 $\phi_{CY,J,t}$，使得 $\omega_{CY,t}=dd^c\phi_{CY,J,t}$，且当 $t\to 0$ 时 $\|\phi_{CY,J,t}-\phi_0\circ\mathrm{Log}_{\mathcal X}\|_{C^0}\to 0$。则在把 $U_{J,t}$ 略微收缩之后，有 $C^\infty$-渐近 $\|\phi_{CY,J,t}-\phi_0\circ\mathrm{Log}_{\mathcal X}\|_{C^k_{loc}}\to 0$。
\end{proposition}

% [Tutor Note: 命题 6.3 原文标为 "(cf. [53, section 4.5])"，是「参见」而非命名，故括注保留 cf. 的含义、写成文献位置。$C^k_{loc}$ 的 loc 是下标（局部），$k$ 是上标，不要抄成 $C^k_{loc}$ 之外的形式；$C^\infty$-渐近指对每个 $k$ 都有 $C^k_{loc}$ 收敛。Alexandrov 解是实 MA 方程的测度论弱解（见 4.5 节）。]

\begin{proof}[证明（概要）]
\begin{itemize}
\item 第一个要素是：由实 MA 方程的正则性理论（参看第 5.4.5 节），在删去 $\mathrm{Int}(\Delta_J)$ 的一个 Hausdorff $(n-1)$-测度为零的子集之后，$\phi_0$ 是光滑的。再把剩下的开集略微收缩，$\phi_0$ 就有 $C^k$ 界。
\item 第二个要素是 Savin 小扰动定理（参看第 5.4.8 节）。过渡到局部万有覆盖后，$\phi_{CY,J,t}$ 与 $\phi_0\circ\mathrm{Log}_{\mathcal X}$ 都解一个复 Monge--Ampère 方程。作为第 5.3.1 节体积形式渐近的推论，它们右端的差在 $t\to 0$ 极限下以任意高的 $C^k$ 范数消没。Savin 的结果于是把 $\phi_{CY,J,t}$ 与 $\phi_0\circ\mathrm{Log}_{\mathcal X}$ 的 $C^0$ 接近性改进为 $C^\infty_{loc}$ 接近性（在 $U_{J,t}$ 小幅收缩之后）。\proofend
\end{itemize}
\end{proof}

% [Tutor Note: 原文 "Hausdorff(n−1)-measure zero" 指 Hausdorff $(n-1)$-测度为零，不是「Hausdorff 维数 $n-1$」。证明标题原文为 "Proof. (Sketch)"，故用可选参数「证明（概要）」；原文末尾有 $\square$，故保留 \proofend。这里用两次「收缩」：先删奇异集、再缩开集，两步不可合并。]

% [原文 p.36]
命题 5.6.3 蕴含在略微收缩的 $U_{J,t}$ 上有半平坦度量渐近，其 $C^\infty$ 误差在 $t\to 0$ 极限下很小：
\begin{equation}
\omega_{CY,t}\sim dd^c(\phi_0\circ\mathrm{Log}_{\mathcal X})=\frac{\sqrt{-1}}{4\pi|\log|t||^2}\sum_{i,j}\frac{\partial^2\phi_0}{\partial x_i\partial x_j}d\log z_i\wedge d\,\overline{\log z_j}. \tag{5.19}
\end{equation}
就 Riemann 度量张量而言，
\begin{equation}
g_{CY,t}\sim\frac{1}{2\pi|\log|t||^2}\mathrm{Re}\Big\{\sum_{1\le i,j\le n}\frac{\partial^2\phi_0}{\partial x_i\partial x_j}d\log z_i\otimes d\log\bar z_j\Big\}. \tag{5.20}
\end{equation}
特别地，Riemann 曲率在 $U_{J,t}$ 中一致有界。由第 5.4.6 节综述的特殊拉格朗日扰动理论，对数映射的 $T^n$ 纤维可以在略微收缩的子集上经 $C^\infty$ 小扰动变成特殊拉格朗日纤维化。

% [Tutor Note: (19) 末项原文的上划线覆盖 "log z_j" 整个（即 $d\overline{\log z_j}$），而 (20) 末项的上划线只覆盖 $z_j$（即 $d\log\bar z_j$）------同一对象、两式排印不一致，按页图分别照抄。另注意 (19) 的分母是 $4\pi|\log|t||^2$、(20) 是 $2\pi|\log|t||^2$（相差因子 2，来自全纯形式与 Riemann 张量的换算），且 (19) 求和只写 $\sum_{i,j}$（无范围），(20) 写 $\sum_{1\le i,j\le n}$。$|\log|t||^2$ 是 $(\log|t|)^2$（外层绝对值再平方）。]

若命题 5.6.3 的假设在 $Sk(X)$ 的所有 $n$ 维面上成立，则所有 $U_{J,t}$ 的并覆盖 $X_t$ 上几乎所有的 CY 测度，而区域收缩过程中损失的测度可忽略。于是 SYZ 猜想的弱度量版本随之成立。

\begin{remark}{5.14}
一个细微之处是：实 Monge--Ampère 方程的局部正则性理论允许 Hausdorff 余维 $1+\epsilon$ 的奇点。这意味着 Kontsevich--Soibelman 猜想中余维二的奇点预言不能简单地从上述论证推出。人们必须在 $Sk(X)$ 上（而不仅在其 $n$ 维面的内部）找到更整体的论证。
\end{remark}

% [Tutor Note: 注 14 是本章的关键限制：「局部」正则性只给出余维 $1+\epsilon$ 的奇异集控制，而 Kontsevich--Soibelman 猜想的预言是「余维二」；两者的差就是为什么必须离开各面的内部、到 $Sk(X)$ 的骨架上作整体论证。$\epsilon$ 的大小不受控。]

\subsection{策略 I：非阿基米德几何}
剩下的任务是：使局部势局部地 $C^0$-收敛到实 MA 方程在 $Sk(X)$ 的开 $n$ 维面上的解（参看命题 5.6.3）。第一个策略 [53] 是：
\begin{itemize}
\item 在 $Sk(X)$ 上解实 MA 方程，不依赖 $X_t$ 上的 CY 度量；
\item 然后试图把该解与 $X_t$ 上 CY 度量的势作比较。首先，需要在 $X_t$ 上构造一个 Kähler 度量，使其局部势在某种拓扑下 $C^0$-接近 $Sk(X)$ 上的实 MA 解。其次，需要某个版本的 $L^1$-体积稳定性估计（参看第 5.4.3 节），以说明该 Kähler 度量与 CY 度量之间的相对势至少在一般区域（generic region）中 $C^0$-小。
\end{itemize}

\subsubsection{NA 几何的动机}
上述策略包含许多问题：

% [原文 p.37]
\begin{itemize}
\item 正如第 5.5.6 节所讨论的，人们不知道如何在 $Sk(X)$ 上直接表述实 MA 方程，也不知道这种表述所需的凸函数类的确切定义。
\item 本质骨架是单纯复形，而 $X_t$ 是复流形。二者在概念上非常不同，我们需要一个拓扑把两边统一起来。
\item 多复变势理论的论证要求 Kähler 势具有整体正性（即 psh 性质）（参看注 5.6）。因此当我们把实 MA 解从 $Sk(X)$ 嫁接到 $X_t$ 上时，必须保证整体正性。困难在于非一般区域（non-generic regions），那里 $X_t$ 的复结构高度奇异。
\end{itemize}

这些问题自然地指向 NA 几何：
\begin{itemize}
\item NA MA 与实 MA 比较性质是产生解的自然途径；
\item 混合拓扑是把 $X_t$ 与含本质骨架的 $X_K^{an}$ 作比较的自然拓扑；
\item 半正度量的概念是 NA 几何内置的。
\end{itemize}

\begin{remark}{5.15}
非阿基米德途径的一个副产品是：Calabi--Yau 局部势的极限实际上与子序列无关，因为非阿基米德版本的 Calabi--Yau 度量已知是唯一的。
\end{remark}

% [Tutor Note: 注 15 回应了 p36 之前的 "for a subsequence" 表述：定理 6.1/6.2 只能沿子序列取极限，而注 15 说明若用 NA 途径，极限其实不依赖子序列（因 NA CY 度量唯一）。"cf. Remark 6" 指向更早章节对 psh 整体正性的讨论，不在本章。]

\subsubsection{嫁接实 MA 解}
设 $(\mathcal X,\mathcal L)$ 是半稳定的 snc 模型，且 $\mathcal L|_X=L$。NA 多复变势理论给出 $X_K^{an}$ 上 $L$ 的连续半正度量 $\|\cdot\|_{CY}=\|\cdot\|_{\mathcal L}e^{-\phi_0}$，它解 NA MA 方程 (5.16)；我们此后假设它满足 NA MA 与实 MA 比较性质，于是 $\phi_0$ 在本质骨架 $Sk(\mathcal X)$ 的各 $n$ 维开面 $\mathrm{Int}(\Delta_J)$ 上解实 MA 方程（参看第 5.5.5 节）。

\begin{proposition}{5.6.4}[53, Lemma 4.1, 4.2]
给定任意 $\epsilon\ll 1$，并令 $t$ 依 $\epsilon$ 足够小。则存在 Kähler 度量 $\omega_{\psi,t}$，使得
\begin{itemize}
\item 在 $\mathrm{Log}_{\mathcal X}^{-1}(\mathrm{Int}(\Delta_J))$ 上，可以选取 $\omega_{\psi,t}$ 的局部 Kähler 势 $\phi_{J,t}$ 满足 $|\phi_{J,t}-\phi_0\circ\mathrm{Log}_{\mathcal X}|<\epsilon$；
\item 总变差 $\int_{X_t}\Big|\frac{|\log|t||^n\omega_{\psi,t}^n}{(L^n)}-d\mu_t\Big|<\epsilon$；
\item $\omega_{\psi,t}$ 相对于某个固定 Fubini--Study 参考度量的 Kähler 势一致有界，且与 $t,\epsilon$ 无关。
\end{itemize}
\end{proposition}

% [Tutor Note: 本页 p37 的 "essential skeleton Sk(𝒳)" 用的是花体 $\mathcal X$（与前面 p32、p34 的 $Sk(X)$ 不同）------原文本身在这两页混用 $Sk(X)$ 与 $Sk(\mathcal X)$，译文按页图分别照抄。命题 6.4 第二条的分母是相交数 $(L^n)$，分子是 $|\log|t||^n\omega_{\psi,t}^n$（$n$ 次幂），整个绝对值符号包住「归一化体积形式减 $d\mu_t$」；这是「总变差」的度量。$\omega_{\psi,t}$ 的下标 $\psi$ 指该度量的 Kähler 势函数 $\psi$。]
